\documentclass[10pt,a4paper]{amsart}
\usepackage[margin=19mm]{geometry}
\usepackage{amsmath,amssymb,mathtools}
\usepackage[scr=rsfs,cal=euler]{mathalfa}
\usepackage{xfrac}
\usepackage[unicode,hidelinks,bookmarks=false]{hyperref}
\newcommand{\ie}{i.e.\ }
\newcommand{\cf}{cf.\ }
\newcommand{\resp}{respectively\ }
\newcommand{\Subsection}{\subsection}
\providecommand{\nobreakdash}{\nobreak\hbox{-}}
\setlength{\emergencystretch}{2em}
\multlinegap0pt
\usepackage{bm}
\usepackage{bbm}
\usepackage{dsfont}
\usepackage{xspace}
\usepackage{cancel}
\usepackage{mathtools}
\usepackage{amsthm}
\makeatletter
\@ifundefined{claim}{\newtheorem*{claim}{Claim}}{}
\@ifundefined{thm}{\newtheorem{thm}{Theorem}}{}
\@ifundefined{cor}{\newtheorem{cor}[thm]{Corollary}}{}
\@ifundefined{lemma}{\newtheorem{lemma}[thm]{Lemma}}{}
\makeatother
\DeclareMathOperator{\sch}{SymConv}
\newcommand{\Pcal}{\mathcal P}
\newcommand{\T}{\mathcal T}
\newcommand{\eps}{\epsilon}
\title[Asymmetry of 2-step Transit Probabilities in 2-Coloured Regu]{Asymmetry of 2-step Transit Probabilities in 2-Coloured Regular Graphs}
\author{Ron Gray, J. Robert Johnson}
\date{Source: arXiv:2307.06054v2 (CC BY 4.0)}
\begin{document}
\maketitle

\noindent\emph{Source and licence.} Asymmetry of 2-step Transit Probabilities in 2-Coloured Regular Graphs. Version arXiv:2307.06054v2 (CC BY 4.0), distributed under the Creative Commons Attribution 4.0 International licence, as recorded in the arXiv licence metadata for this exact version and shown on the abstract page https://arxiv.org/abs/2307.06054v2. Full rights notice: licenses/arxiv-2307.06054-CC-BY-4.0.txt.\par\smallskip

\noindent\textit{Editorial scope.} Counted unit: the Cor whose source label is \texttt{construction} in the article (source0.tex lines 376--385, source label \texttt{construction}), delivered together with the article's own complete original proof of it (source0.tex lines 400--505, one \texttt{proof} environment of 106 lines). No mathematics is written, repaired or re-derived: statement and proof are the article's own text line for line. The proof is detached from the statement in the source: the article states the result at lines 376--385 and prints its proof at lines 400--505, and the proof environment's own header names the statement, so the association is the article's own. Required context retained uncounted: squares (lines 124--129); main:thm (lines 148--150). Every definition, display and label of those lines is retained unchanged. Omitted from this package and not redistributed: 1--123, 130--147, 151--375, 386--399, 506--610. Nothing inside a retained excerpt is omitted or altered: every retained body is delivered exactly as the article's lines are written, so no omission receipt of any kind accompanies this package. Citations: the retained excerpts carry no citation command at all, so no bibliography entry of the article is transcribed and no reference is redistributed. Reference adaptation: none was needed and none was made; every reference inside the delivered unit resolves inside it, so no unresolved reference or citation is produced. Printed slips kept literal: no printed word, symbol, hypothesis or number of the counted statement or of its proof was altered. Layout and class: the article's own class is \texttt{article} with packages including amsfonts, amsmath, amsthm, graphicx, xcolor, enumerate, fullpage; the wrapper is the lane's pinned \texttt{amsart} editorial wrapper at 10pt on A4, which restates the theorem-like environments the retained text needs and adds the article's own macro definitions that the retained text needs (\texttt{\textbackslash Pcal}, \texttt{\textbackslash T}, \texttt{\textbackslash eps}, \texttt{\textbackslash sch}), with extra wrapper packages (bm, bbm, dsfont, xspace, cancel, mathtools). The delivered numbering is the unit's own. Assets: no retained span contains a figure environment, a graphics-inclusion command, a PDF-page inclusion, a table or a TikZ picture; no image file is copied into this package, the package declares no asset dependency and no image file is redistributed with it. Licence: version arXiv:2307.06054v2 of this article is distributed under the Creative Commons Attribution 4.0 International licence, as recorded in the arXiv licence metadata for this exact version and shown on the abstract page https://arxiv.org/abs/2307.06054v2. A full rights notice is delivered at licenses/arxiv-2307.06054-CC-BY-4.0.txt. Characters: every retained excerpt is pure ASCII, so the delivered engine, XeLaTeX with its default Latin Modern text font, reports no missing character.
\begin{lemma}\label{squares}
If $G$ is an $n$-vertex, $d$-regular graph then
\[
P_2(R)=\frac{2}{nd^2}\sum_{v\in R} d^2_R(v)=\frac{2}{nd^2}\sum_{i=1}^d i^2r_i
\]
\end{lemma}

\begin{thm}\label{main:thm}
Let $G$ be a $d$-regular graph with $|V(G)|$ even and let $D_d \subset [0,1]^2$ be the set $\sch\left(\left\{\left(\frac{l}{d},\frac{l^2}{d^2}\right):0 \leq l \leq d\right\}\right)$. We have $P_2(R,B)\in D_d$ for all balanced colourings $(R,B)$ of $G$. 
\end{thm}

\begin{cor}\label{construction}
Let $X_{2m}$ be the symmetric convex hull:
\[
\sch(\left\{\left(\frac{l}{2m},\left(\frac{l}{2m}\right)^2\right) : l=0,m,2m-2,2m-1,2m\right\}
\]
Then
\[
X_{2m}\subseteq\Pcal(\T^m)\subseteq D_{2m}
\]
\end{cor}

\begin{proof}
By Theorem \ref{main:thm} we have
\[
\Pcal(\T^2)\subseteq\sch\left(\left\{\left(0,0\right),\left(\frac{1}{4},\frac{1}{16}\right),\left(\frac{1}{2},\frac{1}{4}\right),\left(\frac{3}{4},\frac{9}{16}\right),\left(1,1\right)\right\}\right).
\]
By the constructions of Corollary \ref{construction} when $m=2$, $r=1,2$ we have
\[
\sch\left(\left\{\left(0,0\right),\left(\frac{1}{2},\frac{1}{4}\right),\left(\frac{3}{4},\frac{9}{16}\right),\left(1,1\right)\right\}\right)\subseteq\Pcal(\T^2).
\]
So to prove the Theorem it suffices to show that if $(\frac{1}{4},y)\in\Pcal(\T^2)$ then $y\geq\frac{1}{8}$. 

Let $(R,B)$ be a balanced colouring of $[k]^2$. Suppose that $v\in B_1$. Let $v=(a,b)$ and suppose, without loss of generality, that the single neighbour of $v$ in $B$ is $w=(a,b+1)$. We will split $B_1$ into types as follows. We say that $v$ has:

\begin{itemize}
\item[Type 3:] if $w\in B_3\cup B_4$
\item[Type 1:] if $(a-1,b+1),(a+1,b+1)\in R$  
\item[Type 2:] otherwise
\end{itemize}
Note that if $w\in B_1$ then $v$ certainly has Type $1$, but if $w\in B_2$ then $v$ could be of Type 1 or Type 2. Let 
\[
t_i=|\{v\in B_1 : \text{$v$ has Type $i$}\}|
\]

We define an auxiliary bipartite graph $H$ with bipartition $(X,Y)$ where
\begin{align*}
X&=\{ v\in B_1 : \text{$v$ has Type $1$ or Type $2$}\}\\
Y&=R_1\cup R_2\cup R_3
\end{align*}

The edges of $H$ are of two types. If $v\in X$ and $r\in Y$ then $v$ and $r$ are adjacent in $H$ if:
\begin{itemize}
\item $v$ has Type $1$ and $v$ and $r$ are adjacent in $T_k^2$,
\item $v=(i,j)$ has Type $2$ and $r\in\{(i,j\pm1), (i\pm1,j),(i\pm1,j\pm1)\}$.
\end{itemize}

\begin{claim}
If $v\in X$ then $\deg_H(v)\geq 2$
\end{claim}

\begin{proof}[Proof of Claim 1.]
If $v$ has Type $1$ then we know $(a-1,b+1),(a+1,b+1)\in R$. Hence both $(a+1,b)$ and $(a-1,b)$ have at least one neighbour in $R$ and one neighbour in $B$ and so are elements of $Y$. It follows that $v$ is adjacent to both $(a+1,b)$ and $(a-1,b)$ in $H$ and so $\deg_H(v)\geq 2$.

If $v$ has Type $2$ then without loss of generality $(a-1,b+1)\in R, (a+1,b+1)\in B$ (if both these vertices were in $B$ then $v$ would have Type $3$). Hence both $(a-1,b+1)$ and $(a-1,b)$ are in $R$ and have at least one neighbour in $R$ and one neighbour in $B$, and so are elements of $Y$. It follows that $v$ is adjacent to both $(a-1,b+1)$ and $(a-1,b)$ in $H$ and so $\deg_H(v)\geq 2$.
\end{proof}

\begin{claim}
If $r\in Y$ then $\deg_H(r)\leq 3$
\end{claim}

\begin{proof}[Proof of Claim 2.]
Suppose that $r=(i,j)\in Y$. We need to show that among the $8$ points $\{(i,j\pm1), (i\pm1,j),(i\pm1,j\pm1)\}$ at most $3$ of them are elements of $X$ which are adjacent to $(i,j)$ in $H$. We will do this by first showing that we cannot have two such elements of $X$ which are adjacent (in $T_k^2$). We then rule out the two remaining $4$ point configurations.  

Suppose, for a contradiction, that both $(i+1,j)$ and $(i+1,j+1)$ are adjacent to $(i,j)$ in $H$. We must have that $(i+1,j)$ and $(i+1,j+1)$ are both in $B_1$ and so have no other neighbours in $B$. Hence both $(i+1,j)$ and $(i+1,j+1)$ are Type $1$ vertices. Now $(i+1,j+1)$ cannot be adjacent to $(i,j)$ in $H$ because of the definition of the edges of $H$ involving type 1 vertices. 

It follows that the only ways $(i,j)$ can have at least $4$ neighbours in $H$ are for these neighbours to be
\begin{itemize}
\item $(i\pm 1,j), (i,j\pm 1)$

or

\item $(i\pm 1,j\pm 1)$
\end{itemize}
In the first case we have $(i,j)\in R_0$. A contradiction since we must have $(i,j)\in Y=R_1\cup R_2\cup R_3$. In the second case, note that $(i+1,j)$ must be in $R$; if it were in $B$ then $(i+1,j+1)$ would have either Type 1 (if $(i+2,j)\in R$) or Type 3 (if $(i+2,j)\in B$) contradicting that $(i+1,j+1)$ is a neighbour of $(i,j)$ in H. Similarly, $(i-1,j),(i,j+1),(i,j-1)\in R$ and so $(i,j)\in R_4$. Again, this contradicts the definition of $Y$. This establishes the second claim.
\end{proof}

Using these two claims to double count the edges of $H$ we get:
\[
2(t_1+t_2)\leq E(H)\leq 3(r_1+r_2+r_3).
\]

Let $n=k^2$. Suppose that $P_2(R)=\frac{1}{4}+\eps$ then by Lemma \ref{squares}:
\[
\frac{1}{4}+\eps=\frac{1}{8n}\left(r_1+4r_2+9r_3+16r_4\right)
\]
and so
\[
2n+8\eps n=r_1+4r_2+9r_3+16r_4.
\]
We also have 
\[
I(R)=r_1+2r_2+3r_3+4r_4.
\]
Hence we obtain
\begin{equation}\label{pr}
2n+8\eps n=4I(R)-(3r_1+4r_2+3r_3)\leq 4I(R)-3(r_1+r_2+r_3)\leq 4I(R)-2(t_1+t_2)
\end{equation}

Now, consider $P_2(B)$. We have that
\[
8nP_2(B)=b_1+4b_2+9b_3+16b_4.
\]
Using the fact that $I(R)=I(B)=b_1+2b+2+3b_3+4b_4$ we can eliminate $b_2$ from this obtaining 
\begin{align*}
8nP_2(B)&=b_1+2\left(I(R)-b_1-3b_3-4b_4\right)+9b_3+16b_4\\
&=2I(R)-b_1+3b_3+8b_4.
\end{align*}
We know that $b_1=t_1+t_2+t_3$. We also know that $3b_3+8b_4\geq 3b_3+4b_4\geq t_3$ because each Type $3$ vertex in $B_3$ is adjacent to a vertex in $B_3$ or $B_4$. Using these bounds in the expression for $8nP_2(B)$ above gives
\begin{align*}
8nP_2(B)&\geq 2I(R)-(t_1+t_2)\\
&\geq n +4\eps n
\end{align*}
by equation (\ref{pr}). Finally, we obtain
\[
P_2(B)\geq\frac{1}{8}+\frac{\eps}{2}
\]
\end{proof}

\end{document}
